\chapter{ESKF的完整推导：传播、观测与重置}

\label{chap:md-09}

第八章已经沿时间顺序走完一次ESKF循环，但把$\Phi_k,Q_{d,k},H,K$和$G_{\mathrm{reset}}$保留为整体映射。现在从同一组坐标和误差约定出发，追踪这些矩阵的来源。每一节先写清真实量、名义量和局部变化，再把分量关系合为矩阵；读到一个非零块时，应能指出它把哪一种旧误差送到了哪一种新误差。

\section{统一约定与推导入口}

\subsection{状态、误差和坐标系保持同一套约定}

沿用第\ref{chap:md-08}章：$R=R_{WB}$把机体系向量送到世界系，$p,v,g_W$在世界系表达，IMU读数及偏置在机体系表达。状态及误差顺序固定为

\begin{equation*}
x=(R,p,v,b_g,b_a),\qquad
\delta x=(\delta\theta,\delta p,\delta v,\delta b_g,\delta b_a).
\end{equation*}

姿态使用右侧误差$R=\hat R\operatorname{Exp}(\delta\theta^\wedge)$，其余分量使用加法。所有Jacobian都在这套局部坐标下计算；更换扰动方向时，矩阵也必须跟着换。

\subsection{一次传播采用什么数值近似}

这里采用第八章的区间起点积分器：旋转用指数映射更新，比力的世界方向在一步起点近似固定，位置使用旧速度和$\tfrac12 h^2$项。连续动力学用于看清误差耦合，实际离散矩阵则要与选定的名义递推一致。连续噪声谱密度、区间平均噪声样本与积分后的离散协方差分别命名，不能仅因都描述噪声便把它们当作同一个矩阵。

\subsection{本章要补出的计算}

\begin{center}\small\begin{tabular}{p{0.420\textwidth}p{0.420\textwidth}}\toprule

第八章的量 & 本章追踪的来源 \\

\midrule

$\Phi_k,Q_{d,k}$ & 真实与名义运动的一阶差、一步离散化和噪声积分 \\

$H$ & 观测映射在名义状态附近的局部微分 \\

$K,P_u$ & 线性高斯观测的局部后验 \\

$G_{\mathrm{reset}}$ & 注入后更换误差参考点的坐标变换 \\

\bottomrule\end{tabular}\end{center}

\section{ESKF：从真实与名义运动得到连续误差系统}

这一节为式\eqref{eq-intuition-error-step}计算传播模型的来源。先把真实运动和名义运动写在同一坐标约定下，比较两者如何分开；连续矩阵只是这些分量关系的汇总。

\subsection{先把两条轨迹写清楚}

由$\omega_m=\omega+b_g+n_g$和$a_m=a+b_a+n_a$，真实运动的形式方程为

\begin{equation}\label{eq-eskf-true-dynamics}
\begin{aligned}
 \dot R&=R(\omega_m-b_g-n_g)^\wedge,\\
 \dot p&=v,\\
 \dot v&=g_W+R(a_m-b_a-n_a),\\
 \dot b_g&=n_{wg},\qquad\dot b_a=n_{wa}.
\end{aligned}
\end{equation}

$n_g,n_a$是测量噪声，$n_{wg},n_{wa}$驱动偏置随机游走。若采用连续白噪声，它们应按积分噪声理解；下一节专门说明这一记法如何落实为有限方差的增量。

记$\hat\omega=\omega_m-\hat b_g$、$\hat a=a_m-\hat b_a$。计算名义轨迹时把未知噪声取为零，得到

\begin{equation*}
\begin{aligned}
 \dot{\hat R}&=\hat R\hat\omega^\wedge,&\dot{\hat p}&=\hat v,\\
 \dot{\hat v}&=g_W+\hat R\hat a,&
 \dot{\hat b}_g&=0,\quad\dot{\hat b}_a=0.
\end{aligned}
\end{equation*}

非线性系统中，将均值代入动力学通常不等于精确求下一状态的期望；这里的名义轨迹是滤波使用的参考。它的离散计算已经由式\eqref{eq-eskf-nominal-step}给出。

\subsection{把真实状态拆成名义值与误差}

本章的$\oplus$明确指

\begin{equation*}
\begin{aligned}
 R&=\hat R\mathop{\mathrm{Exp}}(\delta\theta^\wedge),& p&=\hat p+\delta p,\\
 v&=\hat v+\delta v,& b_g&=\hat b_g+\delta b_g,\\
 && b_a&=\hat b_a+\delta b_a.
\end{aligned}
\end{equation*}

下面丢弃误差与噪声的二阶乘积。推导的关键是先把真实方程写进同一约定，再减去名义方程，而不是直接从姿态矩阵中相减得到一个九维误差。

\subsection{位置和速度误差}

位置直接相减：

\begin{equation*}
\dot{\delta p}=\frac{d}{dt}(p-\hat p)
               =\dot p-\dot{\hat p}=v-\hat v.
\end{equation*}

真实位置以真实速度变化，名义位置以名义速度变化，因此两条轨迹间的距离变化率就是速度差：

\begin{equation*}
\boxed{\dot{\delta p}=\delta v.}
\end{equation*}

速度多了一层关系。它的误差既可能来自加速度计读数和偏置，也可能来自投影读数所用的姿态。先写出真实输入

\begin{equation*}
a_m-b_a-n_a=\hat a-\delta b_a-n_a.
\end{equation*}

再使用$\mathop{\mathrm{Exp}}(\delta\theta^\wedge)\simeq I+\delta\theta^\wedge$。在乘法展开中，$\delta\theta^\wedge\delta b_a$和$\delta\theta^\wedge n_a$都是小量的乘积，线性化时将它们舍去。

速度方程先展开旋转对比力的作用：

\begin{equation*}
\begin{aligned}
 R(a_m-b_a-n_a)
 &=\hat R\mathop{\mathrm{Exp}}(\delta\theta^\wedge)(\hat a-\delta b_a-n_a)\\
 &\simeq\hat R\hat a-\hat R\hat a^\wedge\delta\theta
                 -\hat R\delta b_a-\hat Rn_a.
\end{aligned}
\end{equation*}

这里使用了$\delta\theta\times\hat a=-\hat a\times\delta\theta$。减去名义速度方程后，已知重力抵消，得到

\begin{equation*}
\boxed{\dot{\delta v}
       =-\hat R\hat a^\wedge\delta\theta
        -\hat R\delta b_a-\hat Rn_a.}
\end{equation*}

姿态估计有误，会把比力投到错误的世界方向；速度误差随后积累成位置误差。因此这三个误差不能作为相互独立的标量噪声处理。

\begin{bookexample}{轻微倾斜怎样变成水平速度误差}

令$\hat R=I$、$\hat a=(0,0,g)^\top$，误差为绕$y$轴的小角度$\delta\theta=(0,\varepsilon,0)^\top$。则

\begin{equation*}
-\hat a^\wedge\delta\theta
 =\delta\theta\times\hat a=(g\varepsilon,0,0)^\top.
\end{equation*}

原本接近竖直方向的比力，因为姿态差异而带来一个水平分量。一小段时间$h$后，这个方向的速度误差约增加$g\varepsilon h$。若$\varepsilon=0.01\,\mathrm{rad}$，对应加速度误差约为$0.098\,\mathrm{m/s^2}$。姿态误差即使看起来很小，也会持续影响速度和位置。

\end{bookexample}

\subsection{姿态误差}

记$E=\hat R^\top R=\mathop{\mathrm{Exp}}(\delta\theta^\wedge)$。

这个相对旋转在名义与真实姿态重合时等于$I$。先对乘积求导：

\begin{equation*}
\dot E=\dot{\hat R}^{\top}R+\hat R^\top\dot R.
\end{equation*}

由$\dot{\hat R}=\hat R\hat\omega^\wedge$和叉积矩阵的反对称性，有$\dot{\hat R}^{\top}=-\hat\omega^\wedge\hat R^\top$。真实角速度又等于$\hat\omega-\delta b_g-n_g$。把这两项代回，而不是对旋转矩阵的九个元素独立取误差，就得到

\begin{equation*}
\dot E=-\hat\omega^\wedge E
          +E(\hat\omega-\delta b_g-n_g)^\wedge.
\end{equation*}

使用$E\simeq I+\delta\theta^\wedge$，保留一阶项，得

\begin{equation*}
\dot{\delta\theta}^{\wedge}
 \simeq-\hat\omega^\wedge\delta\theta^\wedge
       +\delta\theta^\wedge\hat\omega^\wedge
       -(\delta b_g+n_g)^\wedge.
\end{equation*}

叉积矩阵的交换子满足 $[u^\wedge,v^\wedge]=(u\times v)^\wedge$，因而

\begin{equation*}
\boxed{\dot{\delta\theta}
       =-\hat\omega^\wedge\delta\theta-\delta b_g-n_g.}
\end{equation*}

其中$-\hat\omega^\wedge\delta\theta$来自误差所依附的局部坐标随名义姿态变化。若改成左侧误差，这一项及相关坐标会随之改变。

可以用两个情形检查它。若名义姿态暂时不转动，$\hat\omega=0$，只剩$\dot{\delta\theta}=-\delta b_g-n_g$。真实偏置比估计大时，名义输入中扣掉的偏置太少，名义角速度偏大，于是"真实相对名义"的误差向负方向积累，负号正好表达这一关系。

若真实与名义角速度相同，只有已有的小姿态差，那么误差向量也会因随名义姿态转动的表达坐标而改变，这由$-\hat\omega^\wedge\delta\theta$记录。

\subsection{偏置误差与整体矩阵}

真实偏置作随机游走，名义偏置在预测时不变，所以

\begin{equation*}
\dot{\delta b}_g=n_{wg},\qquad\dot{\delta b}_a=n_{wa}.
\end{equation*}

将噪声按陀螺测量、加速度测量、陀螺偏置游走、加速度偏置游走的顺序排列：

\begin{equation*}
n=(n_g^\top,n_a^\top,n_{wg}^\top,n_{wa}^\top)^\top\in\mathbb R^{12}.
\end{equation*}

再汇总得到连续线性误差系统

\begin{equation}\label{eq-eskf-linear-system}
\dot{\delta x}=F_c\delta x+G_cn,
\end{equation}

其中每个块都是$3\times3$，按$(\theta,p,v,b_g,b_a)$排列：

\begin{equation*}
F_c=\begin{bmatrix}
 -\hat\omega^\wedge&0&0&-I&0\\
 0&0&I&0&0\\
 -\hat R\hat a^\wedge&0&0&0&-\hat R\\
 0&0&0&0&0\\
 0&0&0&0&0
 \end{bmatrix},\qquad
 G_c=\begin{bmatrix}
 -I&0&0&0\\
 0&0&0&0\\
 0&-\hat R&0&0\\
 0&0&I&0\\
 0&0&0&I
 \end{bmatrix}.
\end{equation*}

$F_c\in\mathbb R^{15\times15}$，$G_c\in\mathbb R^{15\times12}$。

读块矩阵时，每一行对应一个误差变化率，每一列对应一个误差来源。例如第三行乘上$\delta x$就是

\begin{equation*}
(-\hat R\hat a^\wedge)\delta\theta
 +0\delta p+0\delta v+0\delta b_g+(-\hat R)\delta b_a,
\end{equation*}

恰好还原刚才的速度误差方程。第二行的$I$位于速度列，表示位置误差由速度误差驱动。零块也有含义：在当前模型中，位置误差不会直接改变比力测量，所以速度行的位置列为零。 同样，$F_{vb_a}=-\hat R$把机体系偏置误差转换成世界系速度变化率。每个非零块都可以回到对应分量方程检查。

\section{ESKF：把连续误差变成一步传播与协方差}

现在补齐式\eqref{eq-intuition-error-step}中的$\Phi_k,w_k$及式\eqref{eq-intuition-p-predict}中的$Q_{d,k}$。连续$F_c$告诉我们误差变化率，离散$\Phi_k$告诉我们走完一步后的误差，不能把两者直接视为同一个矩阵。

下面区分两条路线：积分连续线性误差系统，或者直接微分选定的离散递推。它们在小时间步的一阶项上相容，高阶项可以不同。实现时应明确选用哪条路线，并配套使用相应的噪声模型。

\subsection{连续白噪声的含义}

连续模型中的白噪声是广义随机过程，不是每个时刻都能取一个有限方差样本的普通函数。令

\begin{equation*}
n=(n_g^\top,n_a^\top,n_{wg}^\top,n_{wa}^\top)^\top,
 \qquad \mathbb E[n(t)n(s)^\top]=Q_c\delta(t-s).
\end{equation*}

$Q_c$表示连续噪声谱密度矩阵。式\eqref{eq-eskf-true-dynamics}是惯性导航常用的形式记法；后面用积分后的噪声增量及其协方差落实计算。名义积分与误差线性化采用小时间步近似，不在这里展开完整的流形随机微分方程理论。

\subsection{状态转移与积分噪声}

线性系统的状态转移矩阵满足

\begin{equation*}
\frac{\partial}{\partial t}\Phi(t,s)=F_c(t)\Phi(t,s),
 \qquad\Phi(s,s)=I.
\end{equation*}

从$t_k$积分到$t_{k+1}$，得到

\begin{equation*}
\delta x_{k+1}=\Phi_k\delta x_k+w_k,
 \qquad\Phi_k=\Phi(t_{k+1},t_k),
\end{equation*}

\begin{equation*}
w_k=\int_{t_k}^{t_{k+1}}\Phi(t_{k+1},\tau)G_c(\tau)n(\tau)\,d\tau.
\end{equation*}

若连续白噪声与旧状态误差独立，则积分噪声零均值，协方差为

\begin{equation}\label{eq-eskf-qd}
Q_{d,k}=\int_{t_k}^{t_{k+1}}
 \Phi(t_{k+1},\tau)G_c(\tau)Q_cG_c(\tau)^\top
 \Phi(t_{k+1},\tau)^\top\,d\tau.
\end{equation}

因此

\begin{equation*}
\boxed{P_{k+1}^-=\Phi_kP_k^+\Phi_k^\top+Q_{d,k}.}
\end{equation*}

第一项传播已有不确定性，第二项加入本时间段的新噪声。若二者相关，交叉项不能省略。

冻结一步内的连续矩阵时，可取$\Phi_k=\exp(F_ch)$并计算对应积分；再作一阶近似才有

\begin{equation*}
\Phi_k\simeq I+F_ch,\qquad Q_{d,k}\simeq G_cQ_cG_c^\top h.
\end{equation*}

这两式是小时间步近似。连续矩阵$F_c$与离散矩阵$\Phi_k$不是同一个对象。

\subsection{与所选离散积分器一致的误差传播}

若直接对式\eqref{eq-eskf-nominal-step}作局部微分，可以得到与该积分器匹配的离散矩阵。令

\begin{equation*}
\phi_k=\hat\omega_kh,\qquad A_k=\mathop{\mathrm{Exp}}(\phi_k^\wedge),
 \qquad B_k=-\hat R_k\hat a_k^\wedge.
\end{equation*}

指数映射的右Jacobian由

\begin{equation*}
\mathop{\mathrm{Exp}}((\phi+\epsilon)^\wedge)
 =\mathop{\mathrm{Exp}}(\phi^\wedge)\mathop{\mathrm{Exp}}((J_r(\phi)\epsilon)^\wedge)
   +O(\|\epsilon\|^2)
\end{equation*}

定义。对$\theta=\|\phi\|$，其矩阵为

\begin{equation*}
J_r(\phi)=I-\frac{1-\cos\theta}{\theta^2}\phi^\wedge
              +\frac{\theta-\sin\theta}{\theta^3}(\phi^\wedge)^2,
 \qquad J_r(0)=I.
\end{equation*}

小角度时用$I-\frac12\phi^\wedge+\frac16(\phi^\wedge)^2+\cdots$计算，避免小量相除的数值问题。对应离散误差转移为

\begin{equation*}
\Phi_k^{\mathrm{step}}=\begin{bmatrix}
 A_k^\top&0&0&-J_r(\phi_k)h&0\\
 \tfrac12B_kh^2&I&Ih&0&-\tfrac12\hat R_kh^2\\
 B_kh&0&I&0&-\hat R_kh\\
 0&0&0&I&0\\
 0&0&0&0&I
 \end{bmatrix}.
\end{equation*}

例如位置更新中的加速度项对姿态求导，正好产生$\frac12B_kh^2$。这与冻结连续矩阵后取矩阵指数在高阶项上不完全相同：前者微分的是选定数值积分器，后者积分的是冻结后的连续线性系统。应明确选用哪一种近似。

\subsection{噪声密度与离散样本方差}

若把区间平均测量噪声记为$\bar n_{g,k},\bar n_{a,k}$，把整步偏置增量记为$w_{bg,k},w_{ba,k}$，采用在步末加入偏置随机游走的离散模型，则

\begin{equation*}
\eta_k=(\bar n_{g,k}^\top,\bar n_{a,k}^\top,
                 w_{bg,k}^\top,w_{ba,k}^\top)^\top,
\end{equation*}

\begin{equation*}
Q_{\eta,k}=\operatorname{diag}
       (Q_gh^{-1},Q_ah^{-1},Q_{wg}h,Q_{wa}h),
\end{equation*}

其中为简洁假设四类噪声互不相关。该离散模型的噪声矩阵为

\begin{equation*}
L_k=\begin{bmatrix}
 -J_r(\phi_k)h&0&0&0\\
 0&-\tfrac12\hat R_kh^2&0&0\\
 0&-\hat R_kh&0&0\\
 0&0&I&0\\
 0&0&0&I
 \end{bmatrix}.
\end{equation*}

于是用$\Phi_k^{\mathrm{step}}$和$Q_{d,k}^{\mathrm{step}}=L_kQ_{\eta,k}L_k^\top$传播。$Q_g,Q_a,Q_{wg},Q_{wa}$在这里都是连续谱密度块，不是传感器单次读数的方差。

这个区间平均噪声模型也是近似。以一维$\dot v=n_a,\dot p=v$为例，精确连续白噪声积分给出

\begin{equation*}
Q_{pv}=q_a\begin{bmatrix}h^3/3&h^2/2\\h^2/2&h\end{bmatrix},
\end{equation*}

而使用一个恒定区间平均噪声并代入$\frac12n_ah^2$，位置方差为$q_ah^3/4$。差别来自区间内噪声的时间结构。需要更精确传播时，应使用式\eqref{eq-eskf-qd}或与更高阶积分器一致的噪声离散模型，不能仅靠改一个时间因子宣称精确。

这就具体实现了第八章"传播已有不确定性，再加入本步噪声"的两项计算。改变数值积分器时，需要重新求它对状态和噪声的微分；协方差递推的形式仍然保持不变。

\section{ESKF：观测微分与卡尔曼更新的来源}

这一节分别计算式\eqref{eq-intuition-observation}中的$H$，并推导式\eqref{eq-intuition-gain}和式\eqref{eq-intuition-joseph}。前者来自几何映射的微分，后两者来自线性高斯模型的条件分布，作用不同。

预测得到一条名义轨迹和它附近的误差分布。现在外部观测到来，我们先问：如果名义状态正确，传感器本应看到什么？实际观测与这个预测之间的差，才是本次需要解释的信息。

\subsection{先定义创新，再线性化误差}

设欧氏观测为

\begin{equation*}
z=h(x)+v_z,\qquad v_z\sim\mathcal N(0,V),
\end{equation*}

其中$V$是观测噪声协方差，与预测误差独立。为了避免与旋转矩阵$R$冲突，本章用$V$表示观测协方差。预测观测和创新为

\begin{equation*}
\hat z=h(\hat x^-),\qquad r=z-\hat z.
\end{equation*}

将$x=\hat x^-\oplus\delta x$代入，得到

\begin{equation*}
h(\hat x^-\oplus\delta x)
 \simeq h(\hat x^-)+H\delta x,
 \qquad
 H=\left.\frac{\partial h(\hat x^-\oplus\delta x)}{\partial\delta x}
   \right|_{\delta x=0}.
\end{equation*}

于是创新模型为

\begin{equation*}
\boxed{r\simeq H\delta x+v_z.}
\end{equation*}

$H$对误差坐标求导，维数为$m\times15$，不是对旋转矩阵九个元素直接求导得到的矩阵。

这里的求导顺序与第\ref{chap:md-06}章一致：先从误差向量$\delta x$生成合法状态$\hat x^-\oplus\delta x$，再通过$h$得到观测，最后对这个复合函数求导。$H$的每一列回答"沿这一种状态误差方向稍微偏一点，预测观测会怎样变"。

一维位置例子中$h(p)=p$，所以$H=1$，关系退化为$r=\delta p+v_z$；后面的路标例子则会把姿态和位置的影响分别算出来。

\subsection{创新协方差与交叉协方差}

创新中包含两种不确定来源：$H\delta x$来自预测状态，$v_z$来自这次测量。独立且零均值时，两者的协方差相加。状态误差与创新之间的共同变化则来自它们共享的$\delta x$：

\begin{equation*}
\mathbb E[\delta x\,r^\top]
 =\mathbb E[\delta x(H\delta x+v_z)^\top]
 =P^-H^\top.
\end{equation*}

在$\delta x\sim\mathcal N(0,P^-)$的近似下，

\begin{equation*}
S=HP^-H^\top+V,
 \qquad \operatorname{Cov}(\delta x,r)=P^-H^\top.
\end{equation*}

因此近似联合分布为

\begin{equation*}
\begin{bmatrix}\delta x\\r\end{bmatrix}
 \sim\mathcal N\left(0,
 \begin{bmatrix}P^-&P^-H^\top\\HP^-&S\end{bmatrix}\right).
\end{equation*}

即使某个状态块在$H$中的列为零，只要它与被观测状态存在交叉协方差，观测仍可能间接修正它。例如路标观测不直接测量陀螺偏置，但多次运动传播会建立姿态与偏置之间的相关性。

\subsection{观测本身在流形上时}

若观测是旋转，不能套用矩阵相减。可以先选相对旋转，再取局部对数，例如

\begin{equation*}
r=\mathop{\mathrm{Log}}(h(\hat x^-)^\top Z)^\vee.
\end{equation*}

随后在一致的局部噪声与误差约定下推导创新线性化。创新定义、$H$和$V$必须描述同一个局部随机向量。本章后面的具体案例使用欧氏观测，使滤波主线不被额外的观测流形参数化打断。

\subsection{从局部后验的平方项推出增益}

暂时假设$P^-$与$V$正定。固定已经收到的创新$r$，由局部高斯先验与似然，误差后验满足

\begin{equation*}
p(\delta x\mid r)\propto
 \exp\left[-\tfrac12\left(
 \delta x^\top(P^-)^{-1}\delta x
 +(r-H\delta x)^\top V^{-1}(r-H\delta x)\right)\right].
\end{equation*}

先展开观测项中关于$\delta x$的二次式：

\begin{equation*}
\begin{aligned}
 &(r-H\delta x)^\top V^{-1}(r-H\delta x)\\
 &\quad=r^\top V^{-1}r
       -2\delta x^\top H^\top V^{-1}r
       +\delta x^\top H^\top V^{-1}H\delta x.
\end{aligned}
\end{equation*}

其中$r$已是给定数据，第一项不随待估误差变化。把最后一项与先验的二次项合并，再收集线性项，得到

\begin{equation*}
\Lambda_u=(P^-)^{-1}+H^\top V^{-1}H,
 \qquad b=H^\top V^{-1}r.
\end{equation*}

配方后得到

\begin{equation*}
P_u=\Lambda_u^{-1},\qquad
 d:=\widehat{\delta x}=\mathbb E[\delta x\mid r]=P_uH^\top V^{-1}r.
\end{equation*}

这一步是在误差坐标中完成高斯配方：令二次函数梯度为零，得到$\Lambda_ud=b$；二次曲率$\Lambda_u$的逆给出后验协方差。为了避免对整个状态维度的矩阵求逆，把结果改写成对观测维度的矩阵$S$求解。

由矩阵求逆恒等式，

\begin{equation*}
P_uH^\top V^{-1}
 =P^-H^\top(HP^-H^\top+V)^{-1}.
\end{equation*}

因此

\begin{equation*}
\boxed{K=P^-H^\top S^{-1},\qquad \widehat{\delta x}=d=Kr,\qquad
 P_u=P^--P^-H^\top S^{-1}HP^-.}
\end{equation*}

增益把观测创新转换为状态误差修正。它同时利用先验置信度、观测置信度和二者之间的相关结构。实际计算通过分解$S$求解线性方程，不显式形成逆矩阵。即使$P^-$只有半正定，只要$S$可逆，增益形式仍可用于相应的高斯条件更新。

\subsection{Joseph形式从误差展开得到}

线性模型中，更新后的中心化误差为

\begin{equation*}
\delta x-d=(I-KH)\delta x-Kv_z.
\end{equation*}

利用旧误差与观测噪声独立，得到

\begin{equation*}
\boxed{P_u=(I-KH)P^-(I-KH)^\top+KVK^\top.}
\end{equation*}

这就是Joseph形式。在精确代数与上述最优$K$下，它与简化式$(I-KH)P^-$等价；数值上它更容易保持对称半正定结构。

此时$d$仍是旧参考点$\hat x^-$附近的误差均值，$P_u$是该旧误差坐标中围绕$d$的后验协方差。它们尚未完成流形上的状态更新。

回到一维例子，旧参考是$10$，更新后误差分布的中心是$d=1.6$，方差为$0.8$。这也可以描述为位置分布以$11.6$为中心，但目前保存的误差坐标仍然以$10$为零点。下一节就是把这两种描述正式接起来。

\section{ESKF：注入之后的重置Jacobian}

这里计算式\eqref{eq-intuition-reset}中的$G_{\mathrm{reset}}$。第八章已经用移动零点的图解释了目的：$P_u$属于旧误差坐标，需要把同一后验改用新参考状态附近的坐标表示。下面只推导这次换元。

\subsection{把修正注入名义状态}

将$d=(d_\theta,d_p,d_v,d_{b_g},d_{b_a})$注入：

\begin{equation*}
\begin{aligned}
 \hat R^+&=\hat R^-\mathop{\mathrm{Exp}}(d_\theta^\wedge),&
 \hat p^+&=\hat p^-+d_p,\\
 \hat v^+&=\hat v^-+d_v,&
 \hat b_g^+&=\hat b_g^-+d_{b_g},\\
 &&\hat b_a^+&=\hat b_a^-+d_{b_a}.
\end{aligned}
\end{equation*}

注入将局部向量变成合法状态。随后要以$\hat x^+$为新的参考，重新定义真实状态的误差。

\subsection{旋转误差坐标为什么改变}

同一真实姿态满足

\begin{equation*}
R=\hat R^-\mathop{\mathrm{Exp}}(\delta\theta_{\mathrm{old}}^\wedge)
  =\hat R^+\mathop{\mathrm{Exp}}(\delta\theta_{\mathrm{new}}^\wedge),
\end{equation*}

因此

\begin{equation*}
\delta\theta_{\mathrm{new}}
 =\mathop{\mathrm{Log}}\left(\mathop{\mathrm{Exp}}(-d_\theta^\wedge)
            \mathop{\mathrm{Exp}}(\delta\theta_{\mathrm{old}}^\wedge)\right)^\vee.
\end{equation*}

令$\delta\theta_{\mathrm{old}}=d_\theta+e_\theta$，围绕后验均值$d_\theta$线性化。

这里$e_\theta$是旧误差坐标中围绕后验均值的剩余偏离，协方差正是$P_u$的相应块。要把它变成新参考点的误差，需要一个局部换元。

第\ref{chap:md-06}章的右Jacobian$J_r(\phi)$回答如下问题：在旋转向量$\phi$上加一小段$\epsilon$，对应在$\mathop{\mathrm{Exp}}(\phi^\wedge)$右侧乘上多大的小旋转？定义关系为

\begin{equation*}
\mathop{\mathrm{Exp}}((\phi+\epsilon)^\wedge)
 =\mathop{\mathrm{Exp}}(\phi^\wedge)\mathop{\mathrm{Exp}}((J_r(\phi)\epsilon)^\wedge)
   +O(\|\epsilon\|^2).
\end{equation*}

两侧描述同一个变化，左侧用旋转向量加法，右侧用已有旋转后的局部小旋转。$J_r(0)=I$；小角度时$J_r(\phi)\simeq I-\tfrac12\phi^\wedge$。闭式表达及小角度展开见第\hyperref[chap:md-09]{ESKF：把连续误差变成一步传播与协方差}节。 由右Jacobian的定义，

\begin{equation*}
\mathop{\mathrm{Exp}}((d_\theta+e_\theta)^\wedge)
 \simeq\mathop{\mathrm{Exp}}(d_\theta^\wedge)
        \mathop{\mathrm{Exp}}((J_r(d_\theta)e_\theta)^\wedge),
\end{equation*}

故

\begin{equation*}
\boxed{\delta\theta_{\mathrm{new}}\simeq J_r(d_\theta)e_\theta.}
\end{equation*}

这也可以直接从乘法消去看出：将右Jacobian展开式代入旧新误差关系，左侧的$\mathop{\mathrm{Exp}}(-d_\theta^\wedge)$与$\mathop{\mathrm{Exp}}(d_\theta^\wedge)$相消，只剩$\mathop{\mathrm{Exp}}((J_r(d_\theta)e_\theta)^\wedge)$。再取局部对数，得到上式。

在小修正下，$J_r(d_\theta)\simeq I-\frac12d_\theta^\wedge$。这个负号由本章采用的右侧误差约定决定。

欧氏误差满足$\delta p_{\mathrm{new}}=\delta p_{\mathrm{old}}-d_p$，速度与偏置类似。因此整体重置Jacobian及协方差为

\begin{equation*}
G_{\mathrm{reset}}=\operatorname{diag}(J_r(d_\theta),I,I,I,I),
\end{equation*}

\begin{equation*}
\boxed{P^+=G_{\mathrm{reset}}P_uG_{\mathrm{reset}}^\top.}
\end{equation*}

重置也会改变姿态与其他状态的交叉协方差，不能只修正左上角的姿态方差块。

例如姿态与位置的后验交叉块为$P_{u,\theta p}$，重置后变为

\begin{equation*}
P^+_{\theta p}=J_r(d_\theta)P_{u,\theta p}.
\end{equation*}

位置误差采用的世界系加法坐标没有旋转，但它与姿态误差的相关关系仍需随姿态误差坐标一起变换。

\subsection{清零的是误差均值的记录}

在线性化模型下，新误差均值回到零，于是滤波器保存$(\hat x^+,P^+)$开始下一轮。"误差清零"不是说真实状态已经完全正确，也不是把协方差置零。剩余不确定性仍由$P^+$描述；非线性重置后出现的高阶均值项被局部近似忽略。

观测更新、注入和重置可按下面的对象链理解：

\begin{equation*}
(\hat x^-,P^-)
 \longrightarrow(d,P_u)
 \longrightarrow\hat x^+
 \longrightarrow(0,P^+).
\end{equation*}

其中第二步中的协方差还属于旧误差坐标，最后一步才搬到新误差坐标。对纯欧氏加法状态，重置Jacobian为$I$，所以普通卡尔曼滤波中这一步往往不显眼。

\section{ESKF：把完整推导用于一次路标观测}

第八章用$h(x)=R^\top(\ell-p)$说明观测模型。现在把它的$H$具体算出，再走一次增益、注入和重置，检查上述工具怎样用于同一个几何观测。

\subsection{已知路标的相对位置观测}

设已知世界系路标$\ell\in\mathbb R^3$，传感器与机体系重合，测到机体系中的三维相对位置：

\begin{equation*}
z=R^\top(\ell-p)+v_z.
\end{equation*}

IMU先按照前面的名义传播和协方差传播得到$(\hat x^-,P^-)$。为简化记号，本节在求导中省略上标$-$，记$\hat z=\hat R^\top(\ell-\hat p)$。代入右侧扰动后，

\begin{equation*}
\begin{aligned}
 h(\hat x\oplus\delta x)
 &=\mathop{\mathrm{Exp}}(-\delta\theta^\wedge)\hat R^\top
                      (\ell-\hat p-\delta p)\\
 &\simeq\hat z+\hat z^\wedge\delta\theta-\hat R^\top\delta p.
\end{aligned}
\end{equation*}

所以按$(\theta,p,v,b_g,b_a)$排列，

\begin{equation*}
\boxed{H=\begin{bmatrix}\hat z^\wedge&-\hat R^\top&0&0&0\end{bmatrix}.}
\end{equation*}

姿态块的正号来自$-\delta\theta\times\hat z=\hat z\times\delta\theta$。如果是相机像素观测，只需继续复合投影函数$\pi$，在有效深度范围内用$D\pi(\hat z)H$得到像素Jacobian，并在像素空间使用对应噪声协方差。

\subsection{一个能看出更新方向的数值片段}

取$\hat R=I$、$\hat p=0$、$\ell=(1,0,0)^\top$，收到$z=(1,0.02,0)^\top$米，则$r=(0,0.02,0)^\top$。为突出相关方向，假设姿态方差为$0.01I$（弧度平方），位置方差为$0.04I$（米平方），两者当前无交叉协方差，观测方差为$V=0.01I$（米平方）。其他状态在本片段中与姿态、位置不相关。

预测路标距离为一米，$y$方向的创新方差为

\begin{equation*}
S_{yy}=0.01+0.04+0.01=0.06\ \mathrm{m}^2.
\end{equation*}

因此本次更新的非零误差均值为

\begin{equation*}
d_{\theta,z}=-\frac{0.01}{0.06}\,0.02
             \approx-0.00333\ \mathrm{rad},
 \qquad
 d_{p,y}=-\frac{0.04}{0.06}\,0.02
             \approx-0.01333\ \mathrm m.
\end{equation*}

路标在机体系中比预测更偏向$+y$，可以由机器人朝$-y$移动或绕$z$轴顺时针转动解释；两种解释按照先验不确定性分担修正。观测还会在位置和姿态之间建立新的相关性，所以之后不能继续保持它们完全独立。

\subsection{完整循环}

一次带路标观测的滤波循环为：

\begin{enumerate}

\item 用每个IMU采样去偏置，传播名义状态；从同一模型计算$\Phi_k,Q_{d,k}$并传播$P$。

\item 外部观测到来时计算$\hat z,r,H,S$，按需要做创新门限检查。

\item 通过线性求解计算$K=P^-H^\top S^{-1}$，得到$d=Kr$和Joseph协方差$P_u$。

\item 将$d$通过右侧旋转更新与欧氏加法注入名义状态。

\item 用$G_{\mathrm{reset}}$变换整个$P_u$，把局部误差均值记录重置为零。

\end{enumerate}

一个路标的一次三维观测不足以唯一确定全部姿态与位置，更不能直接观测所有速度与偏置。多个路标、跨时刻运动和先验共同决定可观性。

若路标未知，也可扩展$x$加入$\ell_1,\ldots,\ell_N$，其误差采用世界系加法。静态地图点的名义传播不变，观测对当前路标的Jacobian为$\hat R^\top$；必须维护机器人与地图之间的交叉协方差，并在路标初始化时考虑由观测与机器人状态带来的相关性。仅把地图点追加到数组而忽略这些相关性，并没有得到完整的EKF-SLAM。

至此，第八章ESKF流程中保留整体写法的各个矩阵，都有了相应的计算来源。下一章将沿用局部误差与协方差的思想，但把参考对象换成起点固定的一段IMU相对运动。

\section{从区域概率回看ESKF}

第八章已经解释预测改变可能状态的分布、观测改变各候选状态的权重。本节把这一认识写回第\ref{chap:md-07}章的区域概率语言，并给出创新门限的条件。

\subsection{从高斯创新到卡方变量}

在观测与预测匹配、局部模型有效、$S$正定的条件下，$r\approx\mathcal N(0,S)$。令$S=LL^\top$，白化创新$y=L^{-1}r$近似服从$\mathcal N(0,I_m)$，所以

\begin{equation*}
d_M^2=r^\top S^{-1}r=y^\top y\approx\chi_m^2.
\end{equation*}

给定期望接收概率$\alpha$，取$\gamma=F_{\chi_m^2}^{-1}(\alpha)$，则接受区域

\begin{equation*}
A_\gamma=\{r:r^\top S^{-1}r<\gamma\}
\end{equation*}

在上述模型下具有约为$\alpha$的概率。这里使用的是创新协方差$S=HP^-H^\top+V$，不能只拿传感器噪声$V$代替，因为预测状态本身也有误差。

\subsection{接回第\ref{chap:md-07}章的原像语言}

把$r$到$d_M^2$的映射记为$q$，接收区域就是$q^{-1}([0,\gamma))$；再沿残差产生映射取原像，就得到底层样本空间中的接收事件。测度论保证这些事件与概率运算有一致定义，高斯近似则给出具体的卡方阈值。

线性化偏差、错误关联、非高斯噪声都会使实际接收率偏离$\alpha$。此外，一旦按门限筛选了样本，其条件分布已被截断，不能继续把筛选后的经验统计无条件等同于原来的高斯模型。门限用于发现不相容观测，不能替代完整的数据关联与噪声建模。

\subsection{观测为什么会改变各个状态的权重}

第\ref{chap:md-07}章将分布理解为给区域分配概率。预测用运动模型把原来的可能状态向前推进，过程噪声使推进结果继续分散；观测更新则检查每个候选状态能否解释新数据。

先用有限个候选状态说明。若先验权重为$w_i$，观测在该候选状态下的似然为$\ell_i$，则更新后的权重为

\begin{equation*}
w_i^+=\frac{\ell_iw_i}{\sum_j\ell_jw_j}.
\end{equation*}

例如两个候选位置先验各占$1/2$，新观测的似然分别为$0.8$与$0.2$，归一化后权重变为$0.8$与$0.2$。观测没有移动真实机器人，而是改变了我们对候选状态的相对相信程度。

对一般状态分布$\mathsf P^-$，同一规则写成

\begin{equation*}
\mathsf P^+(A)
 =\frac{\int_A\ell(x)\,\mathsf P^-(dx)}
        {\int\ell(x)\,\mathsf P^-(dx)},
\end{equation*}

其中分母假设有限且为正。ESKF把名义状态附近的先验与观测作局部高斯、线性近似，才得到前面的卡尔曼公式。

\subsection{从确定性预测到带噪声的预测}

没有过程噪声时，若下一状态为$f_k(x)$，则预测区域$A$的概率为$\mathsf P_{k-1}^+(f_k^{-1}(A))$。这就是第\ref{chap:md-07}章的分布推前。

有噪声时，从同一个$x$出发也可能到达多个区域。记$\mathcal K_k(x,A;u_k)$为给定$x$和输入$u_k$后下一状态落入$A$的条件概率，便有

\begin{equation*}
\mathsf P_k^-(A)=\int\mathcal K_k(x,A;u_k)\,\mathsf P_{k-1}^+(dx).
\end{equation*}

它对所有可能起点的贡献加权求和。转移核是这个条件概率的统一名称；前面已经推导的名义状态与协方差传播，是对这项分布计算的局部近似。

\begin{quote}\textbf{说明 · 参考状态与均值}\par

在欧氏高斯模型中，参考点可以直接取均值；流形上则先选择局部误差，令其均值在一阶模型下为零。它不自动等于按某种全局距离定义的均值。若选定距离$d$，可另研究最小化$\mathbb E[d(q,X)^2]$的Fréchet均值，但其存在性与唯一性需要额外条件，ESKF的计算不依赖于求出这个全局对象。

\end{quote}

\section{把推导重新接回一次计算}

一次滤波循环可按五步检查：用IMU传播名义状态；从同一运动模型传播$P$；观测到来后计算创新及其Jacobian；求出误差修正并注入；把后验协方差变换到新参考点。每一步都要核对量的坐标系、单位和局部误差顺序。具体实现若换用另一种积分器或旋转扰动约定，矩阵数值会变化，而“先构造合法状态附近的局部映射，再传播或比较误差”的方法不变。
